\section{Introduction}\label{sec:intro}

\begin{figure*}[t]
  \centering
  \includegraphics[width=\textwidth]{20_snap_hyre_pipeline_art.png}
  \caption{Conceptual view of \hyre{} for legal RAG. A first model call turns
  the question into a tentative legal frame and a retrieval-shaped passage; the
  corpus retriever selects evidence; a second model call answers the original
  question from the selected legal text. The precise operational diagram is in
  Figure~\ref{fig:method_flow}.}
  \label{fig:pipeline_art}
\end{figure*}

Legal RAG looks, at first glance, like the cleanest case for retrieval-augmented
generation \citep{lewis2020rag}. Legal questions point to statutes, holdings,
regulations, administrative notices, and other textual authorities. If a system
can retrieve the right passage, the answer should become easier; if retrieval
misses, the answer should become harder. That simple story is attractive, but
it is too coarse for actual legal evaluation. A raw question can have poor
lexical overlap with the governing rule. A generated query can retrieve more
gold evidence while also pulling in stronger distractors. A benchmark may
already have high raw retrieval, leaving generated-query methods with little
room to help. A model may retrieve the right passage and still choose the wrong
multiple-choice option. A legal benchmark can also lack official passage qrels,
forcing retrieval exposure to be reported as a proxy rather than as Hit@k.

This paper evaluates one fixed method under those constraints. \hyre{} is a
two-call generated-query method. The first call is asked to do two things:
commit to a tentative answer in the dataset's required answer format, and write
a short neutral legal-reference passage that states the rule, standard,
statutory condition, or legal relationship that would resolve the question. The
retriever embeds that passage rather than the raw question. The second call
then receives the original question and the retrieved corpus evidence; it does
not receive the first call's snap answer as answer evidence. The current
harness mode is \method{snap_hyre}; older logs may use
\method{rag_snap_hyde_2call}, which is now only a legacy alias for the same
two-call family.

The important point is that \hyre{} is not presented here as a universal legal
RAG recipe. The signed matrix says something more useful. On BarExamQA, where
raw retrieval is extremely weak (Hit@5 0.0142), \hyre{} improves answer
accuracy over raw RAG for every signed model cell: +2.4 pp for Llama 3.1 8B,
+4.0 pp for Gemma 4 26B, and +5.1 pp for Llama 3.3 70B. On HousingQA, the
signed Llama 70B \hyre{} row improves over raw RAG by +5.1 pp, but
question-only HyDE improves much more (+11.6 pp over raw RAG). On Legal-Link-EU
the pattern reverses: raw-question retrieval already retrieves a source or
target gold document in most rows, and generated-query methods generally reduce
final accuracy. On MASLegalBench, answer accuracy is high and retrieval metrics
are source-document proxies because official per-question qrels are absent.

The result is a retrieval-first but answer-honest paper. We report final answer
accuracy, Hit@5/MRR@5 where aligned gold ids exist, source-document proxies
where they do not, oracle gold-passage controls, signed caveats, and coverage
status. A row is not promoted merely because it appears in a log or in an older
paper draft; it must pass the citation gate in \claimgate. This matters because
the repository contains older diagnostic-routing and controller-era artifacts,
including CaseHOLD and LegalBench-SCALR rows. Those rows remain provenance, but
the current exact-scored main matrix is BarExamQA, HousingQA, Legal-Link-EU,
and MASLegalBench across \method{groq-llama8b}, \method{or-gemma4-26b}, and
\method{groq-llama70b}.

We make four contributions.

\begin{enumerate}[leftmargin=*]
\item We give a fixed-method definition of \hyre{} as a two-call,
snap-conditioned generated-query method, with \method{rag_snap_hyde_2call}
treated only as a legacy alias.
\item We assemble a current source-gated legal RAG matrix over four legal
benchmarks and three model/provider axes, keeping missing and active cells
visible instead of smoothing them into a complete table.
\item We separate retrieval exposure from answer conversion. The signed rows
show BarExamQA retrieval/answer gains, HousingQA generated-query gains with
HyDE stronger than \hyre{}, Legal-Link-EU retrieval degradation under generated
queries, and MASLegalBench source-proxy saturation.
\item We package the evaluation with oracle controls, caveat labels, and
regenerable paper-local figures/tables so that every promoted number is tied
back to source logs and sign-off notes.
\end{enumerate}

The rest of the paper is organized as follows. Section~\ref{sec:preliminary}
places \hyre{} among RAG, HyDE, generated-query retrieval, and legal NLP
benchmarks. Section~\ref{sec:method} defines the method ladder, the exact
\hyre{} mechanics, caching, and validation gates. Section~\ref{sec:experiment}
describes datasets, models, and the current signed coverage. Section
~\ref{sec:analysis} reports the answer matrix, retrieval exposure, oracle
controls, and caveats. Section~\ref{sec:conclusion} closes with limitations
and next steps.
